\section{Heterogeneity by Branch-Market Sophistication}
\label{sec:heterogeneity}
\label{sec:results:sophistication}
\label{sec:results:heterogeneity}

The disciplinary mechanism requires depositors willing and able to act
on new risk information.  We proxy for depositor financial
sophistication with a bank-level index of the ZIP-code demographics of
a bank's branch footprint---the deposit-weighted share of branches in
above-median-education, above-median-investment-income ZIP codes---and
split the sample at the 75th percentile of the index
(Section~\ref{sec:data:sophistication}).  The proxy is imperfect: it
captures the demographic composition of the banking market, not the
sophistication of the depositors of record, and cannot distinguish
retail-branch banking from commercial-banking activity in a given
ZIP code.  We interpret significant coefficients on this index as
evidence that the disclosure response is amplified in markets where
depositors have observable characteristics associated with financial
literacy, not as a direct sophistication measurement.

The heterogeneity is in the price and profitability channels, not the
quantity channel.  Table~\ref{tab:soph_triple_2022} reports pooled
triple interactions Post $\times$ High CECL $\times$ High
Sophistication.  High-CECL banks in high-sophistication markets paid a
steeper increase in interest expense: the triple is $+0.044$ percentage
points per quarter ($p=0.020$), about 17--18 additional basis points
annually on top of the pooled effect, and ROA falls by an additional 30
basis points annually ($-0.076$ per quarter, $p=0.010$); NIM is
directionally consistent but imprecise ($-0.031$, $p=0.175$).  The
uninsured-quantity triple, by contrast, is statistically
indistinguishable from zero ($-0.022$, SE $0.468$), so we do not claim
that sophisticated markets withdrew more; on the quantity margin the
composition shift documented in Section~\ref{sec:results:time_deposits}
is broad-based.  A natural reading is that in more sophisticated
markets the replacement funding must be raised on terms closer to
market wholesale rates, so the same quantity response is more expensive
for the bank, though we regard the heterogeneity evidence as
suggestive rather than sharp.

The pattern is not an artifact of post-CECL branch structure:
rebuilding the index with June 2019 SOD weights---branch footprints
fixed before the CECL rollout and the pandemic---leaves the price and
profitability triples essentially unchanged.  The two indices correlate
at 0.98 and reclassify fewer than 3 percent of banks.
